\expandafter\def\csname HistoricalNoteText2\endcsname{Es ist zu beachten, dass der ausgezeichnete \textsc{d'Alembert} in seiner eigenen Darstellung dieses Beweises geometrische Betrachtungen benutzt und $X$ als Abszisse, $x$ als Ordinate einer Kurve betrachtet hat (nach der Gewohnheit aller Mathematiker des ersten Teils dieses Jahrhunderts, bei denen der Begriff der Funktion weniger gebräuchlich war). Da jedoch alle seine Schlussfolgerungen, wenn man nur auf ihren wesentlichen Inhalt sieht, auf keinen geometrischen, sondern auf rein analytischen Grundlagen beruhen und die Ausdrücke imaginäre Kurven und imaginäre Ordinaten etwas hart erscheinen und den heutigen Leser leichter befremden können, habe ich hier eine rein analytische Darstellungsform vorgezogen. Diese Anmerkung habe ich hinzugefügt, damit niemand, der den eigentlichen Beweis von \textsc{d'Alembert} mit dieser kurzen Darstellung vergleicht, vermutet, es sei etwas Wesentliches verändert worden.}

\expandafter\def\csname HistoricalNoteText3\endcsname{Bei dieser Gelegenheit bemerke ich nebenbei, dass sehr viele dieser Reihen auf den ersten Blick äußerst stark zu konvergieren scheinen, zum Beispiel größtenteils diejenigen, mit denen der ausgezeichnete \textsc{Euler} im hinteren Teil der \textit{Inst. Calc. Diff.}, Kapitel VI, S. 441--474 die Summe anderer Reihen möglichst genau bestimmt (die übrigen Reihen auf S. 475--478 können nämlich tatsächlich konvergieren). Soweit ich weiß, hat dies bisher niemand bemerkt. Daher wäre sehr zu wünschen, dass deutlich und streng gezeigt würde, warum derartige Reihen, die anfangs sehr schnell, dann allmählich immer langsamer konvergieren und schließlich immer stärker divergieren, dennoch eine nahezu wahre Summe liefern, sofern man nur nicht zu viele Glieder nimmt, und wie weit man eine solche Summe sicher als genau ansehen darf?}

